\documentclass[12pt,letterpaper]{article}

\usepackage[utf8]{inputenc}
\usepackage[T1]{fontenc}
\usepackage{amsmath}
\usepackage{amssymb}
\usepackage{geometry}
\usepackage[hidelinks]{hyperref}
\usepackage{mathtools}
\usepackage[backend=biber,style=numeric,sorting=none]{biblatex}

\geometry{margin=1in}
\addbibresource{references.bib}

\newcommand{\assoc}{\mathrm{assoc}}

\title{Explicit Association Picard Approximation for PC-SAFT}
\author{}
\date{}

\begin{document}
\maketitle

\section{Purpose And Roadmap}

This note derives a fully explicit approximation for the PC-SAFT association
site fractions. The story starts with the exact implicit association Helmholtz
form used in the literature, then isolates why that form creates an internal
solve and an implicit-sensitivity burden. The explicit approximation is then
derived as a full-matrix, seven-step, damped Picard map applied to the same
mass-action equations. The final reduced site fractions are substituted into
the standard association Helmholtz-energy expression.

Throughout this note, \(X_a^\star\) denotes the exact implicit mass-action
solution. The reduced form is defined by a fixed seven-step recursion and then
uses the ordinary notation \(X_a\) for the reduced site fractions and
\(a^\assoc\) for the association Helmholtz contribution evaluated with those
reduced fractions. With this convention, no separate approximation superscript
or Picard subscript is carried through the final thermodynamic expressions. The
resulting derivatives are derivatives of the reduced algebraic model; they
coincide with derivatives of the implicit mass-action solution only in limiting
cases where the fixed recursion satisfies the original mass-action balance
exactly.

\section{Exact Implicit Association Model}

\begin{equation}
    a^\assoc =
    \sum_{i=1}^{N_c} x_i
    \sum_{A_i\in\Gamma_i} n_{iA}
    \left(
        \ln X^{A_i}
        - \frac{X^{A_i}}{2}
        + \frac{1}{2}
    \right).
    \tag{4}
    \label{eq:association_helmholtz_literature}
\end{equation}
\setcounter{equation}{4}

Equation~\eqref{eq:association_helmholtz_literature} is the association
Helmholtz contribution written in the usual site-fraction form
\cite{Chapman1990,Huang1990,Gross2002}. The outer sum is over mixture
components, \(x_i\) is the mole fraction of component \(i\), \(\Gamma_i\) is
the set of association sites on component \(i\), and \(n_{iA}\) is the
multiplicity of site class \(A_i\). For an explicitly expanded site list,
\(n_{iA}=1\) for every listed site. The site fraction \(X^{A_i}\) is the
fraction of molecules of component \(i\) that are not bonded at site \(A_i\).

The component number density is
\begin{equation}
    \rho_j=x_j\rho ,
    \label{eq:component_number_density}
\end{equation}
where \(\rho\) is the total molecular number density. If a molar density \(c\)
is used instead, the corresponding number density is \(\rho=N_Ac\). The exact
implicit site fractions satisfy the mass-action relation
\begin{equation}
    X^{A_i,\star}
    =
    \left[
        1+
        \sum_j \rho_j\sum_{B_j}
        X^{B_j,\star}\Delta^{A_iB_j}
    \right]^{-1}.
    \label{eq:exact_mass_action}
\end{equation}
The inner sum is over all association sites \(B_j\) on component \(j\), and the
outer sum is over all mixture components. Equivalently, for a site-fraction
vector \(\mathbf{X}\), the mass-action residual is
\begin{equation}
    \mathcal{R}^{A_i}(\mathbf{X};T,\rho,\mathbf{x},\mathbf{p})
    =
    X^{A_i}
    \left(
        1+\sum_j\rho_j\sum_{B_j}X^{B_j}\Delta^{A_iB_j}
    \right)-1,
    \label{eq:mass_action_residual}
\end{equation}
and the exact fractions satisfy
\(\boldsymbol{\mathcal{R}}(\mathbf{X}^{\star};T,\rho,\mathbf{x},\mathbf{p})
=\mathbf{0}\).

The association strength between site \(A\) on component \(i\) and site \(B\)
on component \(j\) is
\begin{equation}
    \Delta^{A_iB_j}
    =
    d_{ij}^{3}\,
    g_{ij}^{\mathrm{hs}}(d_{ij})\,
    \kappa^{A_iB_j}
    \left[
        \exp\!\left(\frac{\varepsilon^{A_iB_j}}{k_B T}\right)-1
    \right].
    \label{eq:association_strength_literature}
\end{equation}
This relation is commonly listed as the hydrogen-bond strength equation in
Gross--Sadowski-style PC-SAFT summaries \cite{Gross2001,Gross2002,Baygi2015}.
It collects the pair volume \(d_{ij}^3\), the hard-sphere contact value
\(g_{ij}^{\mathrm{hs}}\), the association volume \(\kappa^{A_iB_j}\), and the
Mayer factor for the association energy. The Boltzmann constant is \(k_B\), and
\(T\) is the absolute temperature.

The Mayer factor for the association square-well interaction is
\begin{equation}
    F^{A_iB_j}
    =
    \exp\!\left(\frac{\varepsilon^{A_iB_j}}{k_B T}\right)-1 ,
    \label{eq:association_mayer_factor}
\end{equation}
where \(\varepsilon^{A_iB_j}\) is the association energy between the two sites
\cite{Chapman1990,Huang1990,Gross2002}. The temperature-dependent segment
diameter used in the PC-SAFT hard-sphere contact term is
\begin{equation}
    d_i(T)
    =
    \sigma_i
    \left[
        1-0.12\exp\!\left(-\frac{3\varepsilon_i}{k_B T}\right)
    \right],
    \label{eq:segment_diameter_temperature}
\end{equation}
where \(\sigma_i\) is the temperature-independent segment diameter and
\(\varepsilon_i\) is the segment dispersion energy. The association pair
diameter is the arithmetic mean
\begin{equation}
    d_{ij}
    =
    \frac{d_i+d_j}{2}.
    \label{eq:association_pair_diameter}
\end{equation}

The hard-sphere contact value for the pair is
\begin{equation}
    g_{ij}^{\mathrm{hs}}(d_{ij})
    =
    \frac{1}{1-\zeta_3}
    +
    \frac{3D_{ij}\zeta_2}{(1-\zeta_3)^2}
    +
    \frac{2D_{ij}^{2}\zeta_2^{2}}{(1-\zeta_3)^3},
    \qquad
    D_{ij}=\frac{d_id_j}{d_i+d_j},
    \label{eq:hard_sphere_contact_value}
\end{equation}
with packing moments
\begin{equation}
    \zeta_n
    =
    \frac{\pi}{6}\rho\sum_i x_i m_i d_i^n,
    \qquad n\in\{0,1,2,3\}.
    \label{eq:packing_moments}
\end{equation}
Here \(m_i\) is the segment number.

For cross-association between unlike associating components, Gross and Sadowski
use the combining rules
\begin{align}
    \varepsilon^{A_iB_j}
    &=
    \frac{1}{2}
    \left(
        \varepsilon^{A_iB_i}
        +
        \varepsilon^{A_jB_j}
    \right),
    \label{eq:association_energy_mixing}
    \\
    \kappa^{A_iB_j}
    &=
    \sqrt{\kappa^{A_iB_i}\kappa^{A_jB_j}}
    \left[
        \frac{\sqrt{\sigma_i\sigma_j}}
        {\frac{1}{2}(\sigma_i+\sigma_j)}
    \right]^3 .
    \label{eq:association_volume_mixing}
\end{align}
The first rule defines the cross association energy from the pure-component
association energies; the second scales the geometric mean of the association
volumes by the segment-size ratio \cite{Gross2002}.

\section{Implicit Solve And Derivative Burden}

The Helmholtz expression in \eqref{eq:association_helmholtz_literature} is
algebraic once the site fractions are known, but the exact site fractions are
not given by an explicit formula in the general mixture case. They are defined
implicitly by the coupled residual equations
\begin{equation}
    \boldsymbol{\mathcal{R}}
    \left(
        \mathbf{X}^{\star};T,\rho,\mathbf{x},\mathbf{p}
    \right)
    =
    \mathbf{0}.
    \label{eq:implicit_site_fraction_system}
\end{equation}
Consequently, the exact association contribution has the nested form
\begin{equation}
    a^\assoc_{\star}(q)
    =
    a^\assoc
    \left(
        \mathbf{X}^{\star}(q),q
    \right),
    \qquad
    q=(T,\rho,\mathbf{x},\mathbf{p}).
    \label{eq:implicit_assoc_nested_form}
\end{equation}
The dependence of \(\mathbf{X}^{\star}\) on the thermodynamic variables and
parameters is therefore obtained by implicit differentiation, not by direct
substitution into a closed algebraic expression. For a scalar variable \(q\),
the first derivative follows from
\begin{equation}
    \mathcal{R}_{\mathbf{X}}
    \frac{\partial \mathbf{X}^{\star}}{\partial q}
    +
    \mathcal{R}_{q}
    =
    \mathbf{0},
    \qquad
    \frac{\partial \mathbf{X}^{\star}}{\partial q}
    =
    -
    \mathcal{R}_{\mathbf{X}}^{-1}
    \mathcal{R}_{q},
    \label{eq:implicit_first_derivative}
\end{equation}
where \(\mathcal{R}_{\mathbf{X}}\) is the site-fraction residual Jacobian. The
second derivative in two directions \(q\) and \(r\) requires differentiating
the same residual relation again:
\begin{equation}
    \mathcal{R}_{\mathbf{X}}
    \frac{\partial^2 \mathbf{X}^{\star}}{\partial q\,\partial r}
    +
    \mathcal{R}_{qr}
    +
    \mathcal{R}_{q\mathbf{X}}
    \frac{\partial \mathbf{X}^{\star}}{\partial r}
    +
    \mathcal{R}_{r\mathbf{X}}
    \frac{\partial \mathbf{X}^{\star}}{\partial q}
    +
    \mathcal{R}_{\mathbf{X}\mathbf{X}}
    \left[
        \frac{\partial \mathbf{X}^{\star}}{\partial q},
        \frac{\partial \mathbf{X}^{\star}}{\partial r}
    \right]
    =
    \mathbf{0}.
    \label{eq:implicit_second_derivative}
\end{equation}
Thus exact Hessian information requires residual Jacobian solves plus second
and mixed derivatives of the residual. Third-order information introduces
additional tensor contractions and repeated differentiated solves.

Direct automatic differentiation through an iterative nonlinear solve does not
remove this structure. If the iterations are differentiated literally, the
derivative corresponds to the finite numerical path taken by the solver, which
can depend on stopping tolerances, damping, iteration count, and convergence
history. If the exact converged solution is desired instead, the derivative
must be expressed through the implicit relations above. This is mathematically
clean, but it is a heavier derivative contract for repeated gradient, Jacobian,
and Hessian evaluations.

For an IPOPT-style equilibrium nonlinear program, those derivative evaluations
are called many times while the optimizer changes temperature, density, phase
amounts, and composition-like variables. A finite explicit association closure
is worth studying because it replaces the inner implicit association solve with
a fixed algebraic graph. The reduced graph has predictable cost, fixed
smoothness wherever its elementary operations are smooth, and gradients and
Hessians that can be obtained by ordinary chain-rule automatic
differentiation. The tradeoff is explicit: the derivatives are exact for the
reduced association model, not for the original implicit model, so the reduced
model must be validated against the implicit association residuals and the
resulting thermodynamic properties.

\section{Explicit Closure Matrix Form}

The remaining derivation flattens the component-site notation \(A_i\) into a
single site index \(a=1,\dots,N_s\). Site \(a\) belongs to component \(i(a)\),
and the site fraction \(X_a^\star\) is the corresponding
\(X^{A_i,\star}\). If a flattened site represents \(\nu_a\) equivalent sites,
define the site weight
\begin{equation}
    w_a=x_{i(a)}\nu_a .
    \label{eq:site_weight}
\end{equation}
For an explicitly expanded site list, \(\nu_a=1\). The site-pair association
strength \(\Delta_{ab}\) is the flattened form of
\(\Delta^{A_iB_j}\) from \eqref{eq:association_strength_literature}.

Define the nonnegative association matrix
\begin{equation}
    K_{ab}(T,\rho,\mathbf{x},\mathbf{p})
    =
    \rho\,w_b\,\Delta_{ab}(T,\rho,\mathbf{x},\mathbf{p}) .
    \label{eq:k_matrix}
\end{equation}
Then \eqref{eq:exact_mass_action} becomes the fixed-point system
\begin{equation}
    X_a^\star = G_a(\mathbf{X}^\star)
    =
    \frac{1}{1+\sum_{b=1}^{N_s}K_{ab}X_b^\star},
    \qquad a=1,\dots,N_s .
    \label{eq:fixed_point_map}
\end{equation}
The exact PC-SAFT model iterates or otherwise solves \eqref{eq:fixed_point_map}
until the mass-action residual is sufficiently small. The explicit
approximation replaces that solve with a finite, predetermined algebraic
composition of the same map.

\section{Bounded Row-Sum Initializer}

For each site row, define the row-sum association strength
\begin{equation}
    S_a = \sum_{b=1}^{N_s}K_{ab}.
    \label{eq:row_sum_strength}
\end{equation}
If every site fraction feeding row \(a\) were approximated by a single scalar
\(\chi_a\), the scalar balance would be
\begin{equation}
    \chi_a = \frac{1}{1+S_a\chi_a}.
    \label{eq:scalar_row_sum_balance}
\end{equation}
Multiplying by the denominator gives
\begin{equation}
    S_a\chi_a^2+\chi_a-1=0 .
    \label{eq:scalar_row_sum_quadratic}
\end{equation}
The physical root is
\begin{equation}
    \chi_a =
    \frac{-1+\sqrt{1+4S_a}}{2S_a},
    \qquad S_a>0 .
    \label{eq:row_sum_root_direct}
\end{equation}
Rationalizing the numerator yields the numerically stable form
\begin{equation}
    X_a^{(0)}
    =
    \frac{2}{1+\sqrt{1+4S_a}},
    \qquad a=1,\dots,N_s ,
    \label{eq:row_sum_initializer}
\end{equation}
with the continuous limit \(X_a^{(0)}=1\) when \(S_a=0\). Because
\(S_a\ge 0\), this initializer satisfies \(0<X_a^{(0)}\le 1\).

\section{Seven-Step Damped Picard Closure}

Starting from \eqref{eq:row_sum_initializer}, the undamped Picard proposal at
step \(m\) is the exact fixed-point map evaluated at the previous explicit
iterate:
\begin{equation}
    \widehat{X}_a^{(m)}
    =
    G_a(\mathbf{X}^{(m-1)})
    =
    \frac{1}{
        1+\sum_{b=1}^{N_s}K_{ab}X_b^{(m-1)}
    },
    \qquad m=1,\dots,7 .
    \label{eq:picard_proposal}
\end{equation}
The finite approximation uses fixed damping \(\omega=1/2\):
\begin{equation}
    X_a^{(m)}
    =
    (1-\omega)X_a^{(m-1)}
    +
    \omega\widehat{X}_a^{(m)},
    \qquad \omega=\frac{1}{2},
    \qquad m=1,\dots,7 .
    \label{eq:fixed_damping_picard}
\end{equation}
Substituting \(\omega=1/2\) gives the site-wise recurrence
\begin{equation}
    X_a^{(m)}
    =
    \frac{1}{2}X_a^{(m-1)}
    +
    \frac{1}{2}
    \frac{1}{
        1+\sum_{b=1}^{N_s}K_{ab}X_b^{(m-1)}
    },
    \qquad m=1,\dots,7 .
    \label{eq:picard7_step}
\end{equation}
Since \(K_{ab}\ge 0\) and \(0<X_b^{(m-1)}\le 1\), the proposal lies in
\((0,1]\), and the convex combination in \eqref{eq:picard7_step} also lies in
\((0,1]\). The finite explicit approximation is therefore bounded by
construction.

\section{Reduced Site Fractions}

The reduced site fractions are defined by assigning the seventh Picard iterate
to the nominal site-fraction notation:
\begin{equation}
    X_a = X_a^{(7)},
    \label{eq:final_picard_site_fraction}
\end{equation}
where \(X_a^{(0)}\) is given by \eqref{eq:row_sum_initializer} and each
\(X_a^{(m)}\), \(m=1,\dots,7\), is given by \eqref{eq:picard7_step}. Equivalently,
\begin{equation}
    \mathbf{X}
    =
    \left(\mathcal{P}_{1/2}\right)^7
    \left(\mathbf{X}^{(0)}\right),
    \label{eq:picard7_map}
\end{equation}
where the operator notation means repeated application of the component-wise
recurrence in \eqref{eq:picard7_step}.

\section{Reduced Energy And Residual Check}

Under the reduced-form convention, the association Helmholtz contribution is
\eqref{eq:association_helmholtz_literature} evaluated with the reduced site
fractions \(X_a=X_a^{(7)}\). This is the same algebraic Helmholtz expression
used for the implicit model; only the site fractions differ. The residual of
the original mass-action system at the reduced site fractions is
\begin{equation}
    r_a
    =
    X_a
    \left(
        1+\sum_{b=1}^{N_s}K_{ab}X_b
    \right)-1 .
    \label{eq:picard_mass_action_residual}
\end{equation}
The reduced form matches the exact association balance only in cases where
\(r_a=0\) for every site \(a\).

\section{Reduced-Model Derivative Meaning}

All quantities in \eqref{eq:association_helmholtz_literature} are explicit
functions of \(T\), \(\rho\), \(\mathbf{x}\), and the parameter vector
\(\mathbf{p}\). Therefore repeated application of the chain rule through
\eqref{eq:row_sum_initializer} and \eqref{eq:picard7_step} gives
\begin{equation}
    \frac{\partial a^\assoc}{\partial q},
    \qquad
    q\in\{T,\rho,x_i,p_\ell\},
    \label{eq:picard_derivative_contract}
\end{equation}
as an exact derivative of the reduced explicit Helmholtz model. These
derivatives are not generally equal to the derivatives obtained by implicit
differentiation of the exact mass-action solution. The two derivative sets
coincide only when \eqref{eq:picard_mass_action_residual} vanishes under the
same state and parameter assumptions.

\section{Evaluation Summary}

The complete mathematical evaluation is:
\begin{enumerate}
    \item form the full \(K_{ab}\) matrix from
    \eqref{eq:k_matrix};
    \item initialize with \eqref{eq:row_sum_initializer};
    \item apply exactly seven damped Picard steps using
    \eqref{eq:picard7_step};
    \item evaluate \(a^\assoc\) with
    \eqref{eq:association_helmholtz_literature};
    \item evaluate mass-action residuals with
    \eqref{eq:picard_mass_action_residual} whenever the approximation is
    compared against the implicit exact association model.
\end{enumerate}

\printbibliography

\end{document}
